\chapter{Pretraining on Unlabeled Data}\label{ch:pretrain}

In \cref{ch:gpt} we assembled a GPT model whose weights are random numbers. It produces
fluent-looking \emph{shapes} of text---token IDs in the right range, at the right rate---but
no meaning. This chapter turns that untrained network into a language model. We will define
what ``good'' means for a next-token predictor (cross-entropy and perplexity), write the
training loop, train a small model on a short story, look at how sampling choices change what
the model says, save and restore every bit of training state, and finally load the weights
OpenAI released for GPT-2 so that the remaining chapters can start from a competent model.

For a systems reader the more interesting half of the chapter comes after the loop works.
Pretraining is the single most expensive thing anyone does with a language model: frontier
runs burn $10^{24}$--$10^{26}$ floating-point operations on tens of thousands of accelerators
for months. We therefore spend the second half on the \emph{economics} of a training step:
where the bytes live (weights, gradients, optimizer moments, activations), how mixed precision
and activation checkpointing change that budget, how to measure efficiency with model FLOPs
utilization (MFU), how to checkpoint without stalling the cluster, and how scaling laws tell
you how many parameters and tokens to buy with a given compute budget. \Cref{ch:distributed}
then spreads this single-device loop across many devices.

\begin{figure}[t]
\centering
\begin{tikzpicture}[
  stage/.style={draw, rounded corners=2pt, minimum width=2.25cm, minimum height=0.95cm,
    align=center, font=\footnotesize, fill=#1},
  stage/.default=fcompute!15,
  num/.style={circle, fill=sysblue, text=white, font=\scriptsize\bfseries, inner sep=1.2pt},
  arr/.style={-{Stealth[length=2mm]}, thick, sysblue!70}]
\node[stage] (a) at (0,0) {Generate \&\\ measure loss};
\node[stage] (b) at (2.75,0) {Train / val\\ losses};
\node[stage] (c) at (5.5,0) {Training\\ loop};
\node[stage] (d) at (8.25,0) {Decoding\\ strategies};
\node[stage] (e) at (11,0) {Save / load\\ state};
\node[stage=keygreen!15] (f) at (11,-1.8) {Load GPT-2\\ weights};
\node[stage=keygreen!15] (g) at (8.25,-1.8) {Warmup, cosine,\\ clipping};
\node[stage=napkin!20] (h) at (5.5,-1.8) {Memory \&\\ precision};
\node[stage=napkin!20] (i) at (2.75,-1.8) {Throughput\\ \& MFU};
\node[stage=napkin!20] (j) at (0,-1.8) {Scaling laws\\ \& stability};
\foreach \x/\n in {a/1,b/2,c/3,d/4,e/5,f/6,g/7,h/8,i/9,j/10}
  \node[num] at (\x.north west) {\n};
\draw[arr] (a) -- (b); \draw[arr] (b) -- (c); \draw[arr] (c) -- (d); \draw[arr] (d) -- (e);
\draw[arr] (e) -- (f); \draw[arr] (f) -- (g); \draw[arr] (g) -- (h); \draw[arr] (h) -- (i);
\draw[arr] (i) -- (j);
\node[font=\scriptsize\itshape, text=srcgray, anchor=west] at (-1.1,0.95)
  {Modeling path (blue) $\rightarrow$ better optimization (green) $\rightarrow$ systems economics (amber)};
\end{tikzpicture}
\caption{Roadmap of this chapter. Steps 1--6 take an untrained GPT to a pretrained one;
step 7 makes optimization robust; steps 8--10 are the systems view of what a training step
costs and how to spend a compute budget.}
\label{fig:pretrain:roadmap}
\end{figure}

% ======================================================================
\section{Evaluating generative text models}\label{sec:pretrain:eval}\covers{LLM 5.1}

Before we can improve a model we need a number that says how bad it currently is. For
classification, accuracy is the obvious choice; for a generative model it is not, because the
model's job is to put probability on \emph{all} plausible continuations, not to hit a single
label. The standard metric is the average negative log-likelihood of held-out text, which is
also exactly the quantity we minimize during training.

\subsection{Generating text with an untrained model}\label{sec:pretrain:untrained}\covers{LLM 5.1.1}

We reuse \code{GPTModel} from \cref{ch:gpt} with one change: the context window is reduced
from 1{,}024 to 256 tokens. Attention cost grows quadratically with the window
(\cref{ch:attention}), and our first training corpus is a short story of about five thousand
tokens, so long windows would only shrink the number of training examples. Dropout can be set
to $0.1$ or $0$; modern LLM pretraining usually uses $0$ because a single pass over a huge
corpus rarely overfits.

Two tiny helpers convert between text and the \code{(batch, seq)} tensors the model expects.
They exist mostly so that every later listing reads naturally.

\begin{lstlisting}[caption={Text $\leftrightarrow$ token-ID helpers and a first, meaningless generation.},label={lst:pretrain:helpers}]
import torch
from llmsys.gpt import GPTModel, GPT_CONFIG_124M
from llmsys.generate import generate_simple

def text_to_ids(text, tok):
    return torch.tensor(tok.encode(text)).unsqueeze(0)      # add batch dim

def ids_to_text(ids, tok):
    return tok.decode(ids.squeeze(0).tolist())               # drop batch dim

cfg = dict(GPT_CONFIG_124M, context_length=256)
torch.manual_seed(123)
model = GPTModel(cfg).eval()
out = generate_simple(model, text_to_ids("Every effort moves you", tok),
                      max_new_tokens=10, context_size=cfg["context_length"])
print(ids_to_text(out, tok))   # -> prompt followed by 10 random-looking word pieces
\end{lstlisting}

With random weights the output distribution over the 50{,}257-token vocabulary is close to
uniform, so the generated continuation is a sequence of unrelated word pieces. Our goal is to
make the probability of the \emph{actual} next token large.

\subsection{The text generation loss: cross-entropy and perplexity}\label{sec:pretrain:loss}\covers{LLM 5.1.2}

Take a batch of inputs $\mat{X}\in\{0,\dots,V-1\}^{B\times T}$ and the targets $\mat{Y}$, which
are the same token sequences shifted left by one position (\cref{ch:text}). The model maps
$\mat{X}$ to logits $\mat{Z}\in\R^{B\times T\times V}$. For each position $(b,t)$ a softmax
turns the logit row into a distribution,
\begin{equation}
p_\theta(v\mid x_{b,\le t}) \;=\; \frac{\exp(z_{b,t,v})}{\sum_{u=1}^{V}\exp(z_{b,t,u})},
\label{eq:pretrain:softmax}
\end{equation}
and training wants the probability assigned to the true next token $y_{b,t}$ to be as close to
one as possible. Multiplying probabilities over positions underflows immediately, so we work
with logarithms and average over the $BT$ positions:
\begin{equation}
\mathcal{L}(\theta) \;=\; -\frac{1}{BT}\sum_{b=1}^{B}\sum_{t=1}^{T}
  \log p_\theta\!\left(y_{b,t}\mid x_{b,\le t}\right)
\;=\; -\frac{1}{BT}\sum_{b,t}\Bigl(z_{b,t,y_{b,t}} - \operatorname{logsumexp}_u z_{b,t,u}\Bigr).
\label{eq:pretrain:ce}
\end{equation}
This is the \emph{cross-entropy} between the one-hot data distribution and the model's
distribution, and the right-hand form is how it should be computed: subtracting the
log-sum-exp directly never forms tiny probabilities that would round to zero.
\Cref{fig:pretrain:lossflow} traces the computation for one sequence.

\begin{figure}[t]
\centering
\begin{tikzpicture}[
  box/.style={draw, rounded corners=1.5pt, minimum height=0.8cm, align=center,
    font=\footnotesize, inner sep=3pt},
  arr/.style={-{Stealth[length=1.8mm]}, thick},
  lab/.style={font=\scriptsize, text=srcgray, align=left}]
\node[box, fill=fdata!15, minimum width=2.4cm] (x) at (0,0) {inputs\\ \texttt{(B, T)} ids};
\node[box, fill=fcompute!15, minimum width=2.4cm] (z) at (3.7,0) {logits $\mat Z$\\ \texttt{(B, T, V)}};
\node[box, fill=fcompute!15, minimum width=2.4cm] (lp) at (7.4,0) {log-softmax\\ \texttt{(B, T, V)}};
\node[box, fill=fcompute!15, minimum width=2.4cm] (g) at (11.1,0) {gather at $y_{b,t}$\\ \texttt{(B, T)}};
\node[box, fill=warnred!12, minimum width=2.4cm] (l) at (11.1,-1.8) {$-$mean\\ scalar $\mathcal L$};
\node[box, fill=napkin!20, minimum width=2.4cm] (ppl) at (7.4,-1.8) {perplexity\\ $e^{\mathcal L}$};
\node[box, fill=fdata!15, minimum width=2.4cm] (y) at (11.1,1.7) {targets $\mat Y$\\ \texttt{(B, T)} ids};
\draw[arr] (x) -- node[above, lab]{GPT} (z);
\draw[arr] (z) -- node[above, lab]{$z-\mathrm{lse}(z)$} (lp);
\draw[arr] (lp) -- (g);
\draw[arr] (y) -- (g);
\draw[arr] (g) -- (l);
\draw[arr] (l) -- (ppl);
\node[lab, anchor=north west] at (-1.3,-1.1) {Largest tensor in the whole step:\\
  $B\!\cdot\!T\!\cdot\!V$ floats. For $B{=}8$, $T{=}1024$,\\ $V{=}50{,}257$ in fp32: 1.65\,GB.};
\end{tikzpicture}
\caption{From token IDs to a scalar loss. The target IDs only select one entry per row of the
log-probability tensor; everything else in that $V$-wide row only enters through the
normalizer.}
\label{fig:pretrain:lossflow}
\end{figure}

\begin{lstlisting}[caption={Cross-entropy from scratch (Lab 7, Task 1). It matches \code{F.cross\_entropy} to floating-point precision.},label={lst:pretrain:ce}]
def cross_entropy_from_scratch(logits, targets):
    logits = logits.float()                                  # (B, T, V)
    lse = torch.logsumexp(logits, dim=-1, keepdim=True)      # (B, T, 1)
    log_probs = logits - lse                                 # log softmax
    target_lp = log_probs.gather(-1, targets.unsqueeze(-1)).squeeze(-1)
    return -target_lp.mean()

# PyTorch's fused version expects (N, V) logits and (N,) targets:
loss = torch.nn.functional.cross_entropy(logits.flatten(0, 1), targets.flatten())
\end{lstlisting}

\paragraph{Perplexity.} The exponential of the loss, $\mathrm{PPL}=e^{\mathcal L}$, has a
pleasant reading: it is the size of a uniform distribution that would be equally surprised by
the data. A freshly initialized GPT-2 assigns roughly uniform probability to all $V=50{,}257$
tokens, so its loss is close to $\ln 50{,}257 \approx 10.82$ and its perplexity close to
$50{,}257$. A model with perplexity 20 is, on average, ``choosing among 20 equally likely
tokens''.

\begin{keyidea}[The initial-loss sanity check]
At initialization the logits are small and nearly equal, so the loss should be within a few
tenths of $\ln V$. If your first loss is 30 or 3, something is broken (logit scaling,
targets not shifted, padding counted as targets, wrong vocabulary size). This one-line check
catches a surprising fraction of pipeline bugs before any GPU-hours are spent.
\end{keyidea}

\begin{pitfall}[Perplexities are tokenizer-dependent]
Loss is measured \emph{per token}. Two models with different tokenizers see different numbers
of tokens for the same text, so their perplexities are not comparable. When comparing across
tokenizers, report bits per byte (or per character): multiply the per-token loss by
tokens/bytes and divide by $\ln 2$.
\end{pitfall}

\begin{systemslens}[The logits tensor is the memory hog of the loss]
Everything before the output head scales with $\dmodel$; the logits scale with the vocabulary.
For GPT-2 small at $B=8$, $T=1024$ the logits hold $8\cdot1024\cdot50{,}257\approx4.1\times10^8$
values: \SI{1.65}{GB} in fp32, plus the same again for their gradient, plus the softmax output
if the implementation materializes it. With Qwen-style vocabularies of 152k tokens it is three
times worse. Production stacks therefore never materialize the full logits during training:
\begin{itemize}
  \item \textbf{Chunked / fused cross-entropy} computes the head matmul, log-sum-exp and
    gradient chunk by chunk over the sequence (e.g.\ the Liger kernels, or \code{torch.compile}
    fusing the head and loss), so only one chunk of logits is alive at a time.
  \item \textbf{Vocabulary-parallel cross-entropy} (Megatron-LM) shards the output head across
    tensor-parallel ranks; each rank computes a partial max and partial sum-of-exponentials over
    its vocabulary slice and two tiny all-reduces combine them (\cref{ch:distributed}).
\end{itemize}
Always cast logits to fp32 before the log-sum-exp; bf16 has only 8 bits of mantissa and the
normalizer of a 50k-way softmax loses most of its precision otherwise.
\end{systemslens}

\subsection{Training and validation losses}\label{sec:pretrain:trainval}\covers{LLM 5.1.3}

Our first corpus is deliberately tiny: Edith Wharton's 1908 short story ``The Verdict'' (public
domain), the same text used to build the tokenizer pipeline in \cref{ch:text}. It is about
20{,}000 characters, which the GPT-2 tokenizer turns into roughly 5{,}100 tokens. That is small
enough to train on a laptop in minutes and large enough to watch every phenomenon of
pretraining appear---including overfitting.

We split the \emph{characters} 90/10 into training and validation text, tokenize each part, and
cut both into non-overlapping windows of the context length using the sliding-window dataset of
\cref{ch:text} with \code{stride = context\_length} (\cref{fig:pretrain:split}). The validation
loader keeps \code{drop\_last=False} and \code{shuffle=False} so that every validation token is
seen and results are deterministic. A useful sanity check is to count the tokens in each
loader: with a context of 256 and batch size 2, the 90\% split yields only nine training
batches and the 10\% split a single validation batch. A corpus this small is a toy, and we
will treat the results as such.

\begin{figure}[t]
\centering
\begin{tikzpicture}[x=0.95cm,y=0.8cm]
\foreach \i in {0,...,8} {
  \draw[fill=fdata!25] (\i*1.2,0) rectangle ++(1.1,0.7);
  \node[font=\tiny] at (\i*1.2+0.55,0.35) {win \the\numexpr\i+1\relax};
}
\draw[fill=warnred!20] (10.9,0) rectangle ++(1.1,0.7);
\node[font=\tiny] at (11.45,0.35) {val};
\draw[decorate,decoration={brace,amplitude=5pt}] (0,0.85) -- (10.7,0.85)
  node[midway,above=5pt,font=\scriptsize] {training text: first 90\% of characters $\rightarrow$ 9 windows of 256 tokens};
\draw[decorate,decoration={brace,amplitude=5pt,mirror}] (10.9,-0.15) -- (12,-0.15)
  node[midway,below=5pt,font=\scriptsize, align=center] {last 10\%};
\draw[-{Stealth}, thick, sysblue] (0.2,-0.4) -- (1.0,-0.4) node[right, font=\scriptsize] {target = input shifted by one token};
\node[font=\scriptsize, text=srcgray, align=left, anchor=west] at (0,-1.3)
 {stride = context length $\Rightarrow$ no overlap, each token used once per epoch; the ragged tail is dropped.};
\end{tikzpicture}
\caption{Splitting a short corpus into training and validation windows.}
\label{fig:pretrain:split}
\end{figure}

The function that evaluates a loader simply averages the per-batch loss over at most
\code{num\_batches} batches (\code{llmsys.train.loss\_loader}); limiting the number of batches
keeps periodic evaluation cheap on large datasets. Evaluated before any training, both losses
come out near $\ln V\approx 11$, as expected.

\begin{napkinmath}[How big is ``small''?]
5{,}100 tokens is nothing. GPT-2 was trained on WebText, about 40\,GB of text (on the order of
$10^{10}$ tokens). Llama~3 used about $1.5\times10^{13}$ tokens. At a chat-like
\num{3000} tokens per day, a person would need roughly 14 million years to read Llama~3's
training set. Our corpus is one afternoon of reading---so the model will memorize it.
\end{napkinmath}

% ======================================================================
\section{Training an LLM}\label{sec:pretrain:train}\covers{LLM 5.2}

The pretraining loop is the ordinary deep-learning loop (\cref{ch:pytorch}) applied to the
next-token loss. \Cref{fig:pretrain:loop} shows its structure: an outer loop over epochs, an
inner loop over batches, and the five-beat rhythm of zero-grad, forward, loss, backward,
step. Periodically we evaluate both losses and print a sample generation so a human can see
whether the model is learning language or noise.

\begin{figure}[t]
\centering
\begin{tikzpicture}[
  st/.style={draw, rounded corners=2pt, fill=#1, minimum width=3.2cm, minimum height=0.62cm,
    font=\footnotesize, align=center},
  st/.default=fcompute!12,
  arr/.style={-{Stealth[length=1.8mm]}, thick}]
\node[st=fdata!15] (ep) at (0,0) {for epoch in 1..E};
\node[st=fdata!15] (bt) at (0,-1) {for (x, y) in train\_loader};
\node[st] (zg) at (0,-2) {optimizer.zero\_grad()};
\node[st] (fw) at (0,-3) {loss = CE(model(x), y)};
\node[st] (bw) at (0,-4) {loss.backward()};
\node[st] (sp) at (0,-5) {optimizer.step()};
\node[st=napkin!20] (ev) at (5,-3) {every $k$ steps:\\ train/val loss on\\ a few batches};
\node[st=keygreen!15] (gen) at (5,-4.7) {print a sample\\ generation};
\draw[arr] (ep) -- (bt); \draw[arr] (bt) -- (zg); \draw[arr] (zg) -- (fw);
\draw[arr] (fw) -- (bw); \draw[arr] (bw) -- (sp);
\draw[arr] (sp.west) -- ++(-0.6,0) |- (bt.west) node[pos=0.25, left, font=\scriptsize]{next batch};
\draw[arr, dashed] (sp.east) -- ++(0.5,0) |- (ev.west);
\draw[arr, dashed] (ev) -- (gen);
\draw[arr] (bt.east) -- ++(1.3,0) |- (ep.east) node[pos=0.25, right, font=\scriptsize]{epoch done};
\node[font=\scriptsize, text=srcgray, align=left, anchor=west] at (2.9,-1.3)
  {model.eval() + no\_grad() for evaluation,\\ then back to model.train()};
\end{tikzpicture}
\caption{The pretraining loop. Solid arrows are executed every step; dashed arrows only
periodically, because evaluation and generation are pure overhead for the optimizer.}
\label{fig:pretrain:loop}
\end{figure}

\begin{lstlisting}[caption={A minimal pretraining loop with periodic evaluation.},label={lst:pretrain:simple}]
from llmsys.train import loss_batch, loss_loader

def train_simple(model, train_loader, val_loader, optimizer, device, num_epochs,
                 eval_freq, eval_iter, start_context, tok):
    history, tokens_seen, step = [], 0, 0
    for epoch in range(num_epochs):
        model.train()
        for x, y in train_loader:
            optimizer.zero_grad()
            loss = loss_batch(model, x, y, device)
            loss.backward()
            optimizer.step()
            tokens_seen += x.numel()
            step += 1
            if step % eval_freq == 0:
                model.eval()
                with torch.no_grad():
                    tr = loss_loader(model, train_loader, device, eval_iter)
                    va = loss_loader(model, val_loader, device, eval_iter)
                model.train()
                history.append((step, tokens_seen, tr, va))
                print(f"ep {epoch+1} step {step}: train {tr:.3f} val {va:.3f}")
        print_sample(model, tok, start_context, device)   # eyeball progress
    return history
\end{lstlisting}

We use AdamW with learning rate $4\times10^{-4}$ and weight decay $0.1$. AdamW
(Loshchilov \& Hutter, 2019) is Adam with \emph{decoupled} weight decay: instead of adding
$\lambda\theta$ to the gradient, where Adam's per-coordinate scaling would distort it, the
parameter is shrunk directly by $\eta\lambda\theta$ after the adaptive update. For LLMs this
matters because the adaptive denominators differ by orders of magnitude between embedding
rows of rare and frequent tokens.

\paragraph{What happens on a tiny corpus.} Training GPT-2-small-shaped models for ten epochs on
the nine training batches of ``The Verdict'' produces the textbook picture of overfitting: the
training loss falls from about 10 to well below 1, while the validation loss drops at first
and then flattens out around 6 and starts creeping up. Generated samples turn from noise into
grammatical English---and then into verbatim passages of the story. \Cref{fig:pretrain:overfit}
shows the same phenomenon reproduced on the CPU of the machine that built this book, with a
smaller model and a different small corpus.

\begin{figure}[t]
\centering
\begin{tikzpicture}
\begin{axis}[width=0.85\linewidth, height=5.4cm, xmin=0, xmax=31, ymin=0, ymax=7.2,
  xlabel={epoch}, ylabel={cross-entropy (nats)}, label style={font=\small},
  tick label style={font=\scriptsize},
  legend style={font=\scriptsize, draw=none, at={(0.98,0.97)}, anchor=north east}]
\addplot[very thick, sysblue, mark=*, mark size=1pt] coordinates {(1,6.324) (2,5.892) (3,5.690) (4,5.475) (5,5.280) (6,5.033) (7,4.791) (8,4.554) (9,4.306) (10,4.050) (11,3.829) (12,3.571) (13,3.308) (14,3.100) (15,2.847) (16,2.610) (17,2.406) (18,2.190) (19,1.974) (20,1.779) (21,1.571) (22,1.438) (23,1.257) (24,1.115) (25,1.005) (26,0.857) (27,0.760) (28,0.662) (29,0.590) (30,0.502)};
\addplot[very thick, warnred, mark=square*, mark size=1pt] coordinates {(1,6.406) (2,6.058) (3,5.945) (4,5.872) (5,5.823) (6,5.711) (7,5.632) (8,5.541) (9,5.456) (10,5.383) (11,5.308) (12,5.283) (13,5.212) (14,5.142) (15,5.052) (16,5.016) (17,5.004) (18,4.982) (19,4.976) (20,4.937) (21,4.970) (22,4.970) (23,5.014) (24,5.032) (25,5.073) (26,4.986) (27,5.081) (28,5.109) (29,5.137) (30,5.155)};
\addplot[dashed, srcgray, domain=0:31, samples=2] {6.931};
\draw[densely dotted, keygreen!70!black, thick] (axis cs:20,0) -- (axis cs:20,7.2);
\node[font=\scriptsize, anchor=south west, keygreen!60!black] at (axis cs:20,0.2) {best val (epoch 20)};
\node[font=\scriptsize, anchor=south west, srcgray] at (axis cs:1,6.95) {$\ln V$ = 6.93 (uniform guess)};
\legend{training loss, validation loss}
\end{axis}
\end{tikzpicture}
\caption{Overfitting on a tiny corpus, measured on CPU in this book's build environment: a
3-layer, 96-wide GPT (vocabulary 1{,}024, context 64) trained for 30 epochs on 18k characters
(12 batches of $8\times64$ tokens per epoch) with AdamW, LR $10^{-3}$. Training loss falls
steadily toward zero while validation loss bottoms out at epoch 20 and then rises: the model
is memorizing.}
\label{fig:pretrain:overfit}
\end{figure}

\begin{keyidea}[Overfitting here is a dataset-size artifact]
A 124M-parameter model has enormous capacity relative to 5k tokens: it can store the text.
Real pretraining runs see each token roughly once (one epoch or less), so the validation and
training losses track each other closely and overfitting in the classical sense is rare. When
data is scarce, multiple epochs still help up to about four repetitions, after which returns
diminish quickly (Muennighoff et al., 2023).
\end{keyidea}

% ======================================================================
\section{Decoding strategies to control randomness}\label{sec:pretrain:decoding}\covers{LLM 5.3}

The trained model outputs a distribution over the next token; \emph{decoding} decides which
token to emit. So far we used greedy decoding (\code{argmax}): deterministic and, for a model
that has memorized its training text, a perfect recipe for regurgitation. Sampling from the
distribution instead adds diversity, and two simple knobs control how adventurous the samples
are. \Cref{ch:scaling} returns to these knobs (and adds top-$p$) when we use sampling
deliberately to make reasoning models more accurate.

\subsection{Temperature scaling}\label{sec:pretrain:temp}\covers{LLM 5.3.1}

Replacing \code{argmax} with \code{torch.multinomial(softmax(z))} samples tokens in proportion
to their probabilities. \emph{Temperature} $\tau>0$ divides the logits before the softmax,
\begin{equation}
p_\tau(v) = \frac{\exp(z_v/\tau)}{\sum_u \exp(z_u/\tau)}.
\end{equation}
$\tau<1$ sharpens the distribution toward the argmax (and $\tau\to0$ recovers greedy
decoding); $\tau>1$ flattens it toward uniform. \Cref{fig:pretrain:temperature} applies three
temperatures to an eight-word toy vocabulary.

\begin{figure}[t]
\centering
\begin{tikzpicture}
\begin{axis}[
  width=\linewidth, height=5.2cm, ybar, bar width=4.5pt, ymin=0, ymax=0.75,
  symbolic x coords={the,model,learns,to,predict,next,token,banana},
  xtick=data, x tick label style={font=\scriptsize}, y tick label style={font=\scriptsize},
  ylabel={probability}, ylabel style={font=\small},
  legend style={font=\scriptsize, at={(0.02,0.97)}, anchor=north west, draw=none, fill=white},
  enlarge x limits=0.08]
\addplot[fill=fcompute!60, draw=sysblue] coordinates {(the,0.016) (model,0.002) (learns,0.211) (to,0.0) (predict,0.702) (next,0.064) (token,0.006) (banana,0.0)};
\addplot[fill=fdata!60, draw=keygreen] coordinates {(the,0.069) (model,0.023) (learns,0.253) (to,0.010) (predict,0.462) (next,0.139) (token,0.042) (banana,0.002)};
\addplot[fill=fmemory!60, draw=napkin] coordinates {(the,0.120) (model,0.077) (learns,0.202) (to,0.056) (predict,0.257) (next,0.159) (token,0.099) (banana,0.029)};
\legend{$\tau=0.5$, $\tau=1$, $\tau=2.5$}
\end{axis}
\end{tikzpicture}
\caption{Temperature reshapes the next-token distribution without changing the ranking of
tokens. Logits used: $(2.1,1.0,3.4,0.2,4.0,2.8,1.6,-1.5)$. At $\tau=2.5$ even ``banana''
gets sampled about 3\% of the time.}
\label{fig:pretrain:temperature}
\end{figure}

\subsection{Top-$k$ sampling}\label{sec:pretrain:topk}\covers{LLM 5.3.2}

High temperatures give diversity but also give the long tail of nonsense tokens a real chance.
\emph{Top-$k$} sampling restricts the choice to the $k$ highest-scoring tokens by setting all
other logits to $-\infty$ before the softmax, so their probability becomes exactly zero and the
remaining mass is renormalized among the survivors (\cref{fig:pretrain:topk}). Combining a
moderate top-$k$ (say 25--50) with a temperature above one is a classic recipe for varied but
mostly coherent text.

\begin{figure}[t]
\centering
\begin{tikzpicture}
\begin{axis}[
  width=0.62\linewidth, height=4.6cm, ybar, bar width=6pt, ymin=0, ymax=0.62,
  symbolic x coords={the,model,learns,to,predict,next,token,banana},
  xtick=data, x tick label style={font=\scriptsize, rotate=35, anchor=east},
  y tick label style={font=\scriptsize}, ylabel={probability}, ylabel style={font=\small},
  legend style={font=\scriptsize, at={(0.02,0.97)}, anchor=north west, draw=none},
  enlarge x limits=0.08]
\addplot[fill=fdata!45, draw=keygreen] coordinates {(the,0.069) (model,0.023) (learns,0.253) (to,0.010) (predict,0.462) (next,0.139) (token,0.042) (banana,0.002)};
\addplot[fill=sysblue!70, draw=sysblue] coordinates {(the,0) (model,0) (learns,0.297) (to,0) (predict,0.541) (next,0.163) (token,0) (banana,0)};
\legend{full softmax, top-$k$ ($k{=}3$)}
\end{axis}
\end{tikzpicture}\hfill
\begin{minipage}[b]{0.34\linewidth}\footnotesize
\textbf{Step 1.} Find the $k$-th largest logit $z_{(k)}$.\\[2pt]
\textbf{Step 2.} Set $z_v\leftarrow-\infty$ for every $z_v<z_{(k)}$.\\[2pt]
\textbf{Step 3.} Softmax (with temperature) over what is left; the five masked words get exactly
zero probability.\\[2pt]
\textbf{Step 4.} Sample with \code{torch.multinomial}.\vspace{6pt}
\end{minipage}
\caption{Top-$k$ filtering keeps the $k$ most likely tokens and renormalizes.}
\label{fig:pretrain:topk}
\end{figure}

\subsection{A generation function with temperature, top-$k$ and EOS}\label{sec:pretrain:gen}\covers{LLM 5.3.3}

\Cref{lst:pretrain:generate} folds both knobs into the generation loop and adds an optional
end-of-sequence ID so generation can stop early. \code{temperature=0} keeps the old greedy
behaviour.

\begin{lstlisting}[caption={Generation with temperature, top-$k$, and early stopping (Lab 7, Task 4).},label={lst:pretrain:generate}]
@torch.no_grad()
def generate_text(model, idx, max_new_tokens, context_size, temperature=0.0,
                  top_k=None, eos_id=None, generator=None):
    model.eval()
    for _ in range(max_new_tokens):
        logits = model(idx[:, -context_size:])[:, -1, :].float()   # last position
        if top_k is not None:
            kth = torch.topk(logits, top_k, dim=-1).values[:, -1:]
            logits = torch.where(logits < kth, float("-inf"), logits)
        if temperature > 0.0:
            probs = torch.softmax(logits / temperature, dim=-1)
            nxt = torch.multinomial(probs, 1, generator=generator)
        else:
            nxt = torch.argmax(logits, dim=-1, keepdim=True)
        if eos_id is not None and bool((nxt == eos_id).all()):
            break
        idx = torch.cat([idx, nxt], dim=1)
    return idx
\end{lstlisting}

\begin{systemslens}[Sampling is cheap, but not free, and must be reproducible]
On a GPU the sampling step runs once per generated token on a $(B, V)$ tensor. Top-$k$ is a
partial selection ($O(V)$ with a radix-select kernel) while top-$p$ needs a sort ($O(V\log V)$);
for large batches and vocabularies of 150k tokens inference engines fuse temperature,
top-$k$/top-$p$ and the multinomial draw into one kernel and avoid sorting by rejection-style
sampling over a pivot. Two other systems properties matter more than speed:
\begin{itemize}
\item \textbf{Reproducibility:} pass an explicit \code{torch.Generator} per request (seeded from
  the request ID) rather than relying on the global RNG, which other requests in the same batch
  would advance in a data-dependent order.
\item \textbf{Batch invariance:} with dynamic batching (\cref{ch:serving}) the same request may
  be computed in different batch shapes, which changes matmul reduction orders and therefore
  low-order bits of the logits. Greedy decoding can then flip at near-ties. Bitwise-reproducible
  serving requires batch-invariant kernels.
\end{itemize}
\end{systemslens}

% ======================================================================
\section{Saving and loading model weights}\label{sec:pretrain:save}\covers{LLM 5.4}

Training runs end, crash, and get pre-empted, so persisting state is part of the loop, not an
afterthought. PyTorch's recommended unit of persistence is the \code{state\_dict}: an ordered
dictionary from parameter names to tensors, which decouples the saved data from the Python
class that produced it.

\begin{lstlisting}[caption={Saving and restoring \emph{all} training state (Lab 7, Task 5).},label={lst:pretrain:ckpt}]
def save_checkpoint(path, model, optimizer, step):
    torch.save({"model": model.state_dict(),
                "optimizer": optimizer.state_dict(),   # Adam moments m, v + step counts
                "step": step,
                "rng_cpu": torch.get_rng_state(),      # dropout masks, data order...
                "rng_cuda": torch.cuda.get_rng_state_all()
                            if torch.cuda.is_available() else None}, path)

def load_checkpoint(path, model, optimizer=None):
    ck = torch.load(path, map_location="cpu", weights_only=False)
    model.load_state_dict(ck["model"])
    if optimizer is not None:
        optimizer.load_state_dict(ck["optimizer"])
    torch.set_rng_state(ck["rng_cpu"])
    return ck["step"]
\end{lstlisting}

For inference only the weights are needed (and \code{model.eval()} afterwards, to turn off
dropout). To \emph{continue training}, the optimizer state is just as important: AdamW keeps a
first-moment and second-moment estimate per parameter, and restarting them from zero produces
a burst of overly large updates (the bias correction treats the run as just started). Lab~7
verifies the strongest form of correctness: training 4 steps, checkpointing, restoring into a
freshly constructed model, and training 4 more must give bit-identical weights to training 8
steps straight.

\begin{napkinmath}[Checkpoint sizes]
GPT-2 small has 124M unique parameters. In fp32 the weights take $124\text{M}\times 4 =
\SI{0.50}{GB}$; adding AdamW's two fp32 moments triples that to about \SI{1.5}{GB}. For a 70B
model trained in mixed precision the full training state is about $16$ bytes per parameter
$\approx \SI{1.1}{TB}$---written every few hundred steps by thousands of ranks.
\end{napkinmath}

\begin{systemslens}[Checkpointing at scale]
\begin{itemize}
\item \textbf{Format.} \code{torch.save} pickles Python objects; loading an untrusted pickle can
  execute code. The \emph{safetensors} format stores a JSON header plus raw little-endian tensor
  bytes, so it is safe, can be \code{mmap}-ed for zero-copy lazy loading, and supports reading
  single tensors. It is the de facto format for released weights (the GPT-2 and Qwen3 loaders in
  \code{llmsys} read it).
\item \textbf{Sharded checkpoints.} With ZeRO/FSDP (\cref{ch:distributed}) no rank owns the full
  state. Each rank writes its own shard plus metadata describing the global layout, so a job
  can be resumed on a different number of GPUs by resharding at load time
  (\code{torch.distributed.checkpoint} does this).
\item \textbf{Asynchronous checkpointing.} Copy state to pinned host memory (fast, over PCIe),
  resume training, and let a background thread persist to storage. The stall shrinks from
  ``write \SI{1}{TB} to a parallel file system'' to ``copy a few GB per GPU to host''.
\item \textbf{How often?} Young's approximation gives the checkpoint interval that minimizes
  expected lost work: $T_{\text{opt}}\approx\sqrt{2\,\delta\,M}$, with $\delta$ the time to write
  a checkpoint and $M$ the mean time between failures of the \emph{whole job}. A 60\,s
  checkpoint and a 5-hour job-level MTBF give $T_{\text{opt}}\approx\sqrt{2\cdot60\cdot18{,}000}\,\text{s}
  \approx 25$ minutes. Asynchronous checkpointing shrinks $\delta$, so it also lets you
  checkpoint more often.
\end{itemize}
\end{systemslens}

% ======================================================================
\section{Loading pretrained weights from OpenAI}\label{sec:pretrain:gpt2}\covers{LLM 5.5}

Training a GPT-2-class model on a web-scale corpus costs far more than a laptop can deliver,
but OpenAI released the weights of all four GPT-2 sizes. Loading them into our own
\code{GPTModel} is a good exercise in how architectures and checkpoint layouts relate, and the
resulting model is the starting point of \cref{ch:classify,ch:instruct}.

The original release is a TensorFlow checkpoint plus \code{hparams.json} (which holds
\code{n\_vocab=50257}, \code{n\_ctx=1024}, \code{n\_embd}, \code{n\_head}, \code{n\_layer}).
After reading it into NumPy, the parameters form a nested dictionary with the token and
position embeddings (\code{wte}, \code{wpe}), the final LayerNorm (\code{g}, \code{b}) and a
list of per-block dictionaries. The same weights are also distributed in Hugging Face's
safetensors layout (\code{h.0.attn.c\_attn.weight}, \ldots), which is what
\code{llmsys.gpt.load\_gpt2\_hf\_state\_dict} reads. Either way, three details have to be
handled; \cref{fig:pretrain:mapping} summarizes the mapping.

\begin{enumerate}
\item \textbf{Transposes.} GPT-2's projections are stored as ``Conv1D'' matrices of shape
  \code{(in, out)}, whereas \code{nn.Linear} stores \code{(out, in)}. Every projection is
  transposed on load.
\item \textbf{Fused QKV.} The query, key and value projections are one matrix \code{c\_attn} of
  shape \code{(d, 3d)} (and a bias of length $3d$). We split it into three equal column blocks.
  This also means the released model \emph{has} QKV biases, so the config must set
  \code{qkv\_bias=True}---unlike the bias-free attention we trained from scratch.
\item \textbf{Weight tying.} GPT-2 reuses the token-embedding matrix as the output head. Our
  \code{GPTModel} keeps a separate \code{out\_head}, which is why a naive parameter count of it
  gives 163M rather than 124M; we copy \code{wte} into both.
\end{enumerate}

\begin{figure}[t]
\centering
\begin{tikzpicture}[
  ck/.style={draw, fill=fmemory!18, font=\scriptsize\ttfamily, minimum width=3.6cm,
     minimum height=0.5cm, anchor=west},
  md/.style={draw, fill=fcompute!15, font=\scriptsize\ttfamily, minimum width=4.4cm,
     minimum height=0.5cm, anchor=west},
  arr/.style={-{Stealth[length=1.6mm]}, thin}]
\node[font=\small\bfseries] at (1.8,0.6) {checkpoint (in, out)};
\node[font=\small\bfseries] at (9.2,0.6) {GPTModel (out, in)};
\node[ck] (wte) at (0,0) {wte (V, d)};
\node[ck] (wpe) at (0,-0.65) {wpe (1024, d)};
\node[ck] (qkv) at (0,-1.6) {attn.c\_attn.w (d, 3d)};
\node[ck] (cp) at (0,-2.5) {attn.c\_proj.w (d, d)};
\node[ck] (fc) at (0,-3.15) {mlp.c\_fc.w (d, 4d)};
\node[ck] (mp) at (0,-3.8) {mlp.c\_proj.w (4d, d)};
\node[ck] (ln) at (0,-4.45) {ln\_1/ln\_2 .g, .b};
\node[md] (te) at (7,0.25) {tok\_emb.weight};
\node[md] (oh) at (7,-0.3) {out\_head.weight (tied)};
\node[md] (pe) at (7,-0.85) {pos\_emb.weight};
\node[md] (q) at (7,-1.35) {att.W\_query.weight};
\node[md] (k) at (7,-1.85) {att.W\_key.weight};
\node[md] (v) at (7,-2.35) {att.W\_value.weight};
\node[md] (o) at (7,-2.85) {att.out\_proj.weight};
\node[md] (f1) at (7,-3.35) {ff.layers[0].weight};
\node[md] (f2) at (7,-3.85) {ff.layers[2].weight};
\node[md] (n) at (7,-4.45) {norm1/norm2 .scale, .shift};
\draw[arr] (wte.east) -- (te.west); \draw[arr] (wte.east) -- (oh.west);
\draw[arr] (wpe.east) -- (pe.west);
\draw[arr] (qkv.east) -- node[above, font=\tiny, pos=0.45]{split 3, $\top$} (q.west);
\draw[arr] (qkv.east) -- (k.west); \draw[arr] (qkv.east) -- (v.west);
\draw[arr] (cp.east) -- node[below, font=\tiny, pos=0.5]{$\top$} (o.west);
\draw[arr] (fc.east) -- node[above, font=\tiny, pos=0.5]{$\top$} (f1.west);
\draw[arr] (mp.east) -- node[below, font=\tiny, pos=0.5]{$\top$} (f2.west);
\draw[arr] (ln.east) -- node[below, font=\tiny]{g$\to$scale, b$\to$shift} (n.west);
\end{tikzpicture}
\caption{Mapping GPT-2 checkpoint tensors (per block, left) onto \code{GPTModel} parameters
(right). $\top$ marks a transpose; biases follow their weights (the QKV bias is split into
three, without transposes).}
\label{fig:pretrain:mapping}
\end{figure}

\begin{table}[t]
\centering\small
\begin{tabular}{lrrrrr}
\toprule
model & params & layers $L$ & $\dmodel$ & heads $H$ & weights (fp32) \\
\midrule
gpt2-small  & 124M  & 12 & 768  & 12 & 0.50\,GB \\
gpt2-medium & 355M  & 24 & 1024 & 16 & 1.42\,GB \\
gpt2-large  & 774M  & 36 & 1280 & 20 & 3.10\,GB \\
gpt2-xl     & 1558M & 48 & 1600 & 25 & 6.23\,GB \\
\bottomrule
\end{tabular}
\caption{The four released GPT-2 sizes share one architecture and differ only in depth,
width and head count (all use $V=50{,}257$, context 1{,}024, $\dhead=64$).}
\label{tab:pretrain:gpt2sizes}
\end{table}

\begin{lstlisting}[caption={Core of the OpenAI-layout loader (Lab 7, Task 9). The Hugging Face variant in \code{llmsys/gpt.py} is identical up to key names.},label={lst:pretrain:load}]
def load_openai_gpt2_params(model, params):
    put = lambda param, arr: param.data.copy_(torch.as_tensor(arr))   # + shape check
    put(model.tok_emb.weight, params["wte"])
    put(model.pos_emb.weight, params["wpe"])
    for blk, p in zip(model.trf_blocks, params["blocks"]):
        qw, kw, vw = np.split(p["attn"]["c_attn"]["w"], 3, axis=-1)   # (d, d) each
        qb, kb, vb = np.split(p["attn"]["c_attn"]["b"], 3, axis=-1)
        put(blk.att.W_query.weight, qw.T); put(blk.att.W_query.bias, qb)
        put(blk.att.W_key.weight,   kw.T); put(blk.att.W_key.bias,   kb)
        put(blk.att.W_value.weight, vw.T); put(blk.att.W_value.bias, vb)
        put(blk.att.out_proj.weight, p["attn"]["c_proj"]["w"].T)
        put(blk.ff.layers[0].weight, p["mlp"]["c_fc"]["w"].T)
        put(blk.ff.layers[2].weight, p["mlp"]["c_proj"]["w"].T)
        # ... biases and the two LayerNorms (g -> scale, b -> shift)
    put(model.final_norm.scale, params["g"]); put(model.final_norm.shift, params["b"])
    put(model.out_head.weight, params["wte"])                         # weight tying
\end{lstlisting}

Once loaded, the model writes coherent English from a prompt such as ``Every effort moves
you'', and its loss on ``The Verdict'' is far below that of anything we can train on the story
itself. A strong check that the mapping is right: the pretrained model's validation loss on a
held-out paragraph should be around 3--4 nats, not $\sim$10. A wrong transpose or a swapped
Q/K block yields a loss near the uniform baseline.

\begin{pitfall}[Silent shape coincidences]
Square matrices ($d\times d$) can be loaded with the wrong orientation without any shape error.
Always assert shapes, and always validate a loaded model with a loss on real text or a
reference output, not just ``it runs''.
\end{pitfall}

% ======================================================================
\section{Bells and whistles for the training loop}\label{sec:pretrain:bells}\covers{LLM App.\ D}

The simple loop of \cref{lst:pretrain:simple} works for a toy, but real pretraining is
fragile: early updates can blow up the network before the adaptive optimizer has calibrated its
statistics, late updates with a large step size keep the model bouncing around a minimum, and a
single unlucky batch can produce a gradient large enough to undo hours of progress. Three
techniques address these, and every serious LLM run uses all three.

\subsection{Learning-rate warmup}\label{sec:pretrain:warmup}\covers{LLM D.1}

At initialization, Adam's second-moment estimate $v$ is built from only a few gradients and is a
noisy, typically too-small denominator; at the same time the gradients of a randomly
initialized transformer can be large. Starting at full learning rate therefore produces
destabilizing updates. \emph{Warmup} ramps the learning rate linearly from a small
\code{initial\_lr} to \code{peak\_lr} over the first $W$ optimizer steps,
\begin{equation}
\eta_t = \eta_{\text{init}} + (\eta_{\text{peak}}-\eta_{\text{init}})\,\frac{t}{W}, \qquad t<W.
\end{equation}
Typical warmups are 0.1\%--20\% of total steps; in long pretraining runs a few thousand steps
is common regardless of total length. Implementation-wise the schedule lives \emph{outside}
the optimizer: before each step, write the desired value into every
\code{optimizer.param\_groups[i]["lr"]}.

\subsection{Cosine decay}\label{sec:pretrain:cosine}\covers{LLM D.2}

After warmup, the learning rate should shrink so that late updates refine rather than
overshoot. Cosine decay follows half a cosine period from the peak to a floor
$\eta_{\min}$ (often 10\% of the peak):
\begin{equation}
\eta_t = \eta_{\min} + \tfrac12(\eta_{\text{peak}}-\eta_{\min})
  \Bigl(1+\cos\bigl(\pi\,\tfrac{t-W}{S-W}\bigr)\Bigr), \qquad W\le t\le S,
\label{eq:pretrain:cosine}
\end{equation}
with $S$ the total number of steps. \Cref{fig:pretrain:schedule} shows the combined schedule.

\begin{figure}[t]
\centering
\begin{tikzpicture}
\begin{axis}[width=0.9\linewidth, height=5cm, xmin=0, xmax=150, ymin=0, ymax=0.0105,
  xlabel={optimizer step}, ylabel={learning rate}, ylabel style={font=\small},
  xlabel style={font=\small}, tick label style={font=\scriptsize},
  scaled y ticks=base 10:3,
  legend style={font=\scriptsize, draw=none, at={(0.98,0.97)}, anchor=north east}]
\addplot[very thick, sysblue, domain=0:20, samples=2] {0.0001 + (0.01-0.0001)*x/20};
\addplot[very thick, sysblue, domain=20:150, samples=100, forget plot]
  {0.001 + 0.5*(0.01-0.001)*(1+cos(deg(pi*(x-20)/130)))};
\addplot[thick, dashed, napkin, domain=20:150, samples=2] {0.01};
\addplot[thick, dotted, keygreen, domain=20:150, samples=200]
  {(x<120) ? 0.01 : 0.01 - (0.01-0.001)*(x-120)/30};
\draw[densely dotted, srcgray] (axis cs:20,0) -- (axis cs:20,0.0105);
\node[font=\scriptsize, anchor=west] at (axis cs:21,0.0012) {warmup ends ($W{=}20$)};
\legend{warmup + cosine, constant after warmup, warmup--stable--decay (WSD)}
\end{axis}
\end{tikzpicture}
\caption{Learning-rate schedules for a 150-step run: linear warmup from $10^{-4}$ to
$10^{-2}$, then cosine decay to $\eta_{\min}=10^{-3}$. The dotted curve is the
warmup--stable--decay schedule popular in recent pretraining because a run can be extended
from any point of its stable phase.}
\label{fig:pretrain:schedule}
\end{figure}

\begin{deepdive}[Why cosine, and what replaced it]
Cosine has no deep optimality property; it is a smooth way to spend most of the budget at a
high learning rate and anneal at the end. Its practical weakness is that the total number of
steps $S$ must be fixed in advance---you cannot cheaply continue a run, because the learning
rate has already decayed. \emph{Warmup--stable--decay} (WSD) schedules hold the peak learning
rate for most of training and decay linearly or with a $1-\sqrt{\cdot}$ shape over the last
10--20\% of steps; checkpoints from the stable phase can be branched and annealed at any time,
which makes scaling-law studies far cheaper (Hu et al., 2024, MiniCPM; H\"agele et al., 2024).
\end{deepdive}

\subsection{Gradient clipping}\label{sec:pretrain:clip}\covers{LLM D.3}

Rare batches produce gradients far larger than typical (an unusual document, a numerical
outlier in attention logits). \emph{Global-norm clipping} bounds the size of every update by
rescaling the whole gradient vector $g=(g_1,\dots,g_P)$ when its $\ell_2$ norm exceeds a
threshold $c$ (commonly $c=1$):
\begin{equation}
\|g\|_2=\sqrt{\textstyle\sum_{i} g_i^2}, \qquad
g \leftarrow g\cdot\min\!\Bigl(1, \frac{c}{\|g\|_2}\Bigr).
\end{equation}
Because every coordinate is multiplied by the same factor, the \emph{direction} of the update
is preserved and only its length changes (\cref{fig:pretrain:clip}). PyTorch provides
\code{torch.nn.utils.clip\_grad\_norm\_}; Lab~7 implements it by hand. Call it after
\code{backward()} and before \code{optimizer.step()}---and, when using fp16 loss scaling,
after unscaling the gradients.

\begin{figure}[t]
\centering
\begin{tikzpicture}[scale=1.0]
\draw[->, srcgray] (-0.3,0) -- (4.2,0) node[right, font=\scriptsize] {$g_1$};
\draw[->, srcgray] (0,-0.3) -- (0,3.2) node[above, font=\scriptsize] {$g_2$};
\draw[dashed, keygreen, thick] (1.5,0) arc (0:90:1.5);
\node[font=\scriptsize, keygreen!60!black, anchor=west] at (1.55,0.2) {$\|g\|=c$};
\draw[-{Stealth}, very thick, warnred] (0,0) -- (3.6,2.4) node[right, font=\scriptsize] {raw $g$, $\|g\|=4.3$};
\draw[-{Stealth}, very thick, sysblue] (0,0) -- ({1.5*3.6/4.327},{1.5*2.4/4.327}) node[below right, font=\scriptsize] {clipped};
\draw[-{Stealth}, very thick, keygreen!70!black] (0,0) -- (0.35,1.05) node[above, font=\scriptsize] {small $g$: unchanged};
\begin{scope}[xshift=6.2cm]
\begin{axis}[width=6.4cm, height=4.4cm, xlabel={$\|g\|$ before clipping}, ylabel={$\|g\|$ after},
   xmin=0, xmax=5, ymin=0, ymax=2.2, tick label style={font=\scriptsize},
   label style={font=\scriptsize}]
\addplot[very thick, sysblue, domain=0:5, samples=100] {min(x,1)};
\addplot[dashed, srcgray, domain=0:2.2, samples=2] {x};
\end{axis}
\end{scope}
\end{tikzpicture}
\caption{Global-norm clipping with $c=1$: gradients inside the ball are untouched; larger
ones are scaled back onto it, keeping their direction.}
\label{fig:pretrain:clip}
\end{figure}

\begin{pitfall}[Clipping hides problems as well as fixing them]
Log the pre-clip norm every step (Lab~7's loop records it). A healthy run clips occasionally;
a run that clips on every step is effectively running at a lower, norm-dependent learning rate,
and a steadily rising gradient norm is often the first sign of an impending loss spike.
\end{pitfall}

\subsection{The modified training function}\label{sec:pretrain:trainfn}\covers{LLM D.4}

\Cref{lst:pretrain:trainmodel} combines everything: warmup + cosine, gradient clipping,
gradient accumulation, optional bf16 autocast (\cref{sec:pretrain:mixed}), periodic evaluation
and checkpoints. Note the ordering inside a step: compute the scheduled learning rate, clip,
step, zero the gradients.

\begin{lstlisting}[caption={The full training function from Lab 7 (abridged).},label={lst:pretrain:trainmodel}]
def train_model(model, train_loader, val_loader, optimizer, cfg, start_step=0):
    total = cfg.max_steps or cfg.epochs * (len(train_loader) // cfg.grad_accum)
    warmup = int(cfg.warmup_frac * total)
    step, micro = start_step, 0
    for _ in range(cfg.epochs):
        for x, y in train_loader:
            with torch.autocast("cuda", dtype=cfg.amp_dtype,
                                enabled=cfg.amp_dtype is not None):
                logits = model(x.to(cfg.device))
            loss = F.cross_entropy(logits.float().flatten(0, 1),
                                   y.to(cfg.device).flatten()) / cfg.grad_accum
            loss.backward()                      # grads accumulate in .grad
            micro += 1
            if micro % cfg.grad_accum:
                continue                         # not yet a full (virtual) batch
            lr = warmup_cosine_lr(step, total, cfg.peak_lr, warmup,
                                  cfg.initial_lr, cfg.min_lr)
            for g in optimizer.param_groups:
                g["lr"] = lr
            grad_norm = clip_grad_norm_(model.parameters(), cfg.grad_clip)
            optimizer.step()
            optimizer.zero_grad(set_to_none=True)
            step += 1
            # ... periodic evaluation, logging of lr/grad_norm, checkpointing
\end{lstlisting}

\Cref{fig:pretrain:schedules} compares a constant learning rate with warmup + cosine +
clipping on the same model and data (measured on this book's build machine). The honest
result is that on a small, short run the bells and whistles do not help---and cosine decay can
even hurt, because it anneals before the model has converged. Their value grows with model
size, learning rate and run length, which is precisely where an unstable run is most expensive.
Do not judge these techniques by toy runs; judge them by the failures they prevent at scale.

\begin{figure}[t]
\centering
\begin{tikzpicture}
\begin{groupplot}[group style={group size=2 by 1, horizontal sep=1.4cm},
  width=0.5\linewidth, height=5cm, xmin=0, xmax=210, ymin=4.6, ymax=6.4,
  xlabel={optimizer step}, label style={font=\small}, tick label style={font=\scriptsize},
  legend style={font=\scriptsize, draw=none, at={(0.98,0.97)}, anchor=north east}]
\nextgroupplot[ylabel={validation loss}, title={\footnotesize peak LR $3\times10^{-3}$}]
\addplot[very thick, napkin, mark=*, mark size=1pt] coordinates {(17,6.054) (34,5.946) (51,5.875) (68,5.675) (85,5.540) (102,5.416) (119,5.259) (136,5.115) (153,4.936) (170,4.865) (187,4.818) (204,4.795)};
\addplot[very thick, sysblue, mark=square*, mark size=1pt] coordinates {(17,6.257) (34,5.999) (51,5.982) (68,5.867) (85,5.804) (102,5.696) (119,5.590) (136,5.506) (153,5.406) (170,5.341) (187,5.315) (204,5.297)};
\legend{constant, warmup+cosine+clip}
\nextgroupplot[title={\footnotesize peak LR $1.5\times10^{-2}$}]
\addplot[very thick, napkin, mark=*, mark size=1pt] coordinates {(17,6.075) (34,5.949) (51,5.848) (68,5.643) (85,5.563) (102,5.400) (119,5.245) (136,5.067) (153,4.895) (170,4.881) (187,4.838) (204,4.855)};
\addplot[very thick, sysblue, mark=square*, mark size=1pt] coordinates {(17,6.158) (34,6.043) (51,5.916) (68,5.791) (85,5.752) (102,5.611) (119,5.433) (136,5.271) (153,5.076) (170,4.960) (187,4.896) (204,4.870)};
\legend{constant, warmup+cosine+clip}
\end{groupplot}
\end{tikzpicture}
\caption{Constant learning rate versus warmup (25 steps) + cosine decay (to 10\% of peak) +
clipping at 1.0, measured on CPU in this book's build environment (3-layer GPT, 204 optimizer
steps on 135k characters). At this tiny scale the plain schedule is as good or better: with a
budget of 200 steps, cosine decay lowers the learning rate before the model has finished
learning, and nothing is large enough to blow up. The value of warmup and clipping appears with
large models, high learning rates and long runs, where a single divergence costs days.}
\label{fig:pretrain:schedules}
\end{figure}

% ======================================================================
\section{Where the memory goes}\label{sec:pretrain:memory}

We now switch to the systems view. A training step must keep four kinds of state in device
memory, and only one of them depends on the batch.

\paragraph{Model states.} With $N$ parameters and AdamW in standard mixed precision (bf16 compute,
fp32 master copy), the per-parameter bill is
\begin{center}
\begin{tabular}{lr}
bf16 weights used in forward/backward & 2 bytes \\
bf16 gradients & 2 bytes \\
fp32 master weights & 4 bytes \\
fp32 Adam first moment $m$ & 4 bytes \\
fp32 Adam second moment $v$ & 4 bytes \\
\midrule
total & \textbf{16 bytes}
\end{tabular}
\end{center}
Pure fp32 training without a master copy also costs $4+4+4+4=16$ bytes per parameter. A 7B
model therefore needs $\approx\SI{112}{GB}$ of model states before a single activation is
stored---more than one 80\,GB GPU. This number motivates ZeRO and FSDP in
\cref{ch:distributed}, which shard these 16 bytes across data-parallel ranks.

\paragraph{Activations.} Backpropagation needs intermediate values from the forward pass.
Korthikanti et al.\ (2022) count, for a GPT block with sequence length $s$, micro-batch $b$,
hidden size $h$ and $a$ heads, about
\begin{equation}
M_{\text{act}}^{\text{layer}} \approx s\,b\,h\left(34 + \frac{5\,a\,s}{h}\right)\ \text{bytes (16-bit activations)}.
\label{eq:pretrain:act}
\end{equation}
The 34 collects the inputs of the QKV, output and MLP projections, the LayerNorm inputs, the
GELU input, the dropout masks, and so on; the $5as/h$ term is the $a\times s\times s$ attention
score and probability matrices (plus a dropout mask). FlashAttention (\cref{ch:attention})
never stores those matrices, removing the quadratic term entirely.

\paragraph{Temporary buffers and fragmentation.} Workspace for matmuls and communication
buffers, plus fragmentation in the caching allocator (\cref{ch:pytorch}), typically add
another 5--15\%.

\begin{figure}[t]
\centering
\begin{tikzpicture}
\begin{axis}[xbar stacked, width=0.92\linewidth, height=5.6cm, bar width=10pt,
  xlabel={GiB per replica}, xlabel style={font=\small},
  symbolic y coords={7B-ckpt,7B-flash,7B-naive,1.5B-ckpt,1.5B-flash,1.5B-naive,124M-flash,124M-naive},
  ytick=data, y tick label style={font=\scriptsize}, x tick label style={font=\scriptsize},
  legend style={font=\scriptsize, at={(0.98,0.97)}, anchor=north east, draw=none},
  xmin=0, xmax=210]
\addplot[fill=fmemory!70, draw=napkin] coordinates
 {(1.85,124M-naive) (1.85,124M-flash) (23.2,1.5B-naive) (23.2,1.5B-flash) (23.2,1.5B-ckpt)
  (99.8,7B-naive) (99.8,7B-flash) (99.8,7B-ckpt)};
\addplot[fill=fcompute!60, draw=sysblue] coordinates
 {(8.0,124M-naive) (2.4,124M-flash) (66.8,1.5B-naive) (19.9,1.5B-flash) (1.2,1.5B-ckpt)
  (97.0,7B-naive) (17.0,7B-flash) (1.0,7B-ckpt)};
\draw[dashed, warnred, thick] (axis cs:80,7B-ckpt) -- (axis cs:80,124M-naive);
\node[font=\scriptsize, warnred, anchor=south west, rotate=90] at (axis cs:83,1.5B-ckpt) {one 80 GB GPU};
\legend{model states (16 B/param), activations}
\end{axis}
\end{tikzpicture}
\caption{Memory per data-parallel replica from \cref{eq:pretrain:act} and the 16-byte rule.
124M and 1.5B: $s=1024$, $b=8$; 7B ($L{=}32$, $h{=}4096$, $a{=}32$): $s=4096$, $b=1$.
``naive'' materializes attention scores, ``flash'' does not, ``ckpt'' recomputes each block
and stores only block inputs. Even with checkpointing, the 7B model's states alone exceed one
80\,GB GPU.}
\label{fig:pretrain:membar}
\end{figure}

\begin{napkinmath}[GPT-2 small vs.\ a 7B model]
GPT-2 small ($N=124$M, $L=12$, $h=768$, $a=12$, $s=1024$, $b=8$): model states
$16\times124\text{M}=\SI{1.85}{GiB}$; activations without FlashAttention
$12\cdot 1024\cdot 8\cdot 768\,(34+5\cdot12\cdot1024/768)\approx\SI{8.0}{GiB}$, with FlashAttention
\SI{2.4}{GiB}. At this scale \emph{activations} dominate.

A 7B model ($L=32$, $h=4096$, $a=32$) at $s=4096$, $b=1$: states \SI{99.8}{GiB},
activations \SI{97}{GiB} (naive) or \SI{17}{GiB} (Flash). At this scale \emph{model states}
dominate and no single GPU suffices even at micro-batch 1.
\end{napkinmath}

% ======================================================================
\section{Mixed precision}\label{sec:pretrain:mixed}

Tensor cores deliver their advertised throughput only for 16-bit (or 8-bit) inputs: an H100
SXM does about 989 dense bf16 TFLOP/s versus roughly 67 TFLOP/s for plain fp32 on the CUDA
cores. 16-bit activations also halve activation memory and memory traffic. \emph{Mixed
precision} runs the matmuls in 16 bits while keeping numerically sensitive pieces in fp32.

\begin{table}[t]
\centering\small
\begin{tabular}{lcccl}
\toprule
format & sign/exp/mantissa & max & smallest normal & typical use \\
\midrule
fp32 & 1/8/23 & $3.4\times10^{38}$ & $1.2\times10^{-38}$ & master weights, Adam states \\
bf16 & 1/8/7 & $3.4\times10^{38}$ & $1.2\times10^{-38}$ & activations, weights, grads \\
fp16 & 1/5/10 & 65{,}504 & $6.1\times10^{-5}$ & older GPUs; needs loss scaling \\
fp8 E4M3 & 1/4/3 & 448 & $1.6\times10^{-2}$ & forward GEMM inputs (H100+) \\
fp8 E5M2 & 1/5/2 & 57{,}344 & $6.1\times10^{-5}$ & gradients (H100+) \\
\bottomrule
\end{tabular}
\caption{Floating-point formats in LLM training. bf16 trades mantissa bits for fp32's exponent
range, which is why it rarely needs loss scaling.}
\label{tab:pretrain:formats}
\end{table}

The recipe (Micikevicius et al., 2018) has three ingredients:
\begin{enumerate}
\item \textbf{Autocast.} \code{torch.autocast(device, dtype=torch.bfloat16)} runs matmuls and
  convolutions in bf16 but keeps reductions such as softmax, LayerNorm statistics and the
  loss in fp32. Parameters stay fp32; autocast casts them on the fly (and caches the casts
  within the forward pass).
\item \textbf{fp32 master weights.} The update $\theta\leftarrow\theta-\eta\,\Delta$ is applied
  in fp32. In bf16, with 8 bits of precision, an update smaller than $2^{-8}\approx0.4\%$ of a
  weight rounds away entirely; late in training most updates are that small.
\item \textbf{Loss scaling (fp16 only).} fp16's smallest normal number is $6\times10^{-5}$, and
  many gradient values are smaller. Multiplying the loss by a scale $S$ (e.g.\ $2^{16}$) shifts
  gradients into range; they are divided by $S$ before clipping and the optimizer step.
  \code{torch.amp.GradScaler} adjusts $S$ dynamically, halving it when an overflow (inf/NaN)
  is detected and skipping that step. bf16 has fp32's exponent range and normally needs no
  scaling, which is why it became the default on A100-class and newer hardware.
\end{enumerate}

\begin{systemslens}[FP8 training]
Hopper and later GPUs double tensor-core throughput again with 8-bit floats. FP8 training
(e.g.\ NVIDIA Transformer Engine, DeepSeek-V3's recipe) keeps master weights and optimizer
states in higher precision, casts GEMM inputs to FP8 with \emph{per-tensor or per-block scaling
factors} chosen from recent maxima (``delayed scaling'' or fine-grained tile-wise scales), and
accumulates in higher precision. The engineering problem is outliers: a single huge activation
forces a large scale and flushes everything else in that tensor to zero, which is why
fine-grained (e.g.\ $1\times128$ or $128\times128$ block) scaling is used for the most
sensitive tensors.
\end{systemslens}

% ======================================================================
\section{Trading compute for memory, and measuring efficiency}\label{sec:pretrain:efficiency}

\paragraph{Gradient accumulation.} If the batch that optimization wants (say $10^6$ tokens) does
not fit, split it into $k$ micro-batches, call \code{backward()} on each loss divided by $k$,
and step once. Gradients add in \code{.grad}, so the result equals the large-batch gradient
exactly (when the loss is a mean over equal-sized micro-batches and there is no batch-dependent
layer such as BatchNorm). Lab~7 tests this equivalence. Accumulation costs no extra memory
beyond one micro-batch of activations; its cost is that $k$ forward/backward passes run
sequentially.

\paragraph{Activation checkpointing.} Rather than storing every intermediate, store only the
input of each transformer block and recompute the block's internals during the backward pass.
With $L$ layers and per-layer activation size $A$, memory falls from $LA$ to roughly
$L\cdot a_{\text{in}} + A$ (block inputs plus one block's worth of live intermediates) at the
cost of one extra forward pass---about 33\% more FLOPs, since a backward pass costs about twice
a forward. More generally, checkpointing every $k$-th layer stores $L/k$ checkpoints and
recomputes segments of $k$ layers, so memory is $O(L/k + k)$, minimized at $k=\sqrt L$ for
$O(\sqrt L)$ memory (Chen et al., 2016). \emph{Selective} recomputation (Korthikanti et al.,
2022) recomputes only the cheap-to-recompute but memory-hungry parts (attention softmax and
dropout), recovering most of the memory savings for a few percent of extra compute.

\begin{lstlisting}[caption={Activation checkpointing by wrapping each block (Lab 7, Task 6).},label={lst:pretrain:actckpt}]
from torch.utils.checkpoint import checkpoint

def enable_activation_checkpointing(model):
    for blk in model.trf_blocks:
        inner = blk.forward
        def wrapped(x, cache=None, pad_mask=None, _inner=inner, _blk=blk):
            if _blk.training and cache is None and torch.is_grad_enabled():
                # store only x; recompute attention + MLP internals in backward
                return checkpoint(_inner, x, None, pad_mask, use_reentrant=False)
            return _inner(x, cache=cache, pad_mask=pad_mask)
        blk.forward = wrapped
\end{lstlisting}

\paragraph{Throughput and MFU.} The forward pass of a dense transformer costs about $2N$
FLOPs per token (one multiply and one add per parameter), and the backward pass twice that,
so training costs about
\begin{equation}
C \approx 6\,N\,D \quad\text{FLOPs for $D$ training tokens}
\label{eq:pretrain:6nd}
\end{equation}
(plus an attention term $\approx 12\,L\,h\,s$ per token, which matters only for long contexts).
\emph{Model FLOPs utilization} compares achieved model FLOPs with the hardware peak:
\begin{equation}
\mathrm{MFU} = \frac{6\,N\cdot\text{tokens/s}}{\text{peak FLOP/s}}.
\end{equation}
Recomputation from activation checkpointing does \emph{not} count as model FLOPs (that would be
\emph{hardware} FLOPs utilization, HFU), so MFU honestly penalizes it. Well-tuned large dense
runs report MFUs of roughly 35--55\%; small models on single GPUs are often far lower because
kernels are short and launch overheads and memory-bound elementwise ops dominate.

\begin{napkinmath}[What does 30\% MFU mean for GPT-2 small?]
Suppose GPT-2 small trains at a hypothetical 400{,}000 tokens/s on one H100. Model FLOP/s
$=6\times124\text{M}\times4\times10^5\approx2.98\times10^{14}$, i.e.\ 30\% of the 989 TFLOP/s
bf16 peak. A compute-optimal run for this model ($\approx20$ tokens per parameter, see below)
is $2.5\times10^9$ tokens $\Rightarrow 6\cdot124\text{M}\cdot2.5\text{B}\approx1.8\times10^{18}$
FLOPs, about 1.3 H100-hours at 40\% MFU. GPT-2 small is cheap; data loading and evaluation
overheads, not FLOPs, are what make naive reproductions slow.
\end{napkinmath}

\begin{systemslens}[The usual suspects when MFU is low]
\begin{enumerate}
\item \textbf{Input pipeline stalls.} Tokenize offline into binary shards and memory-map them
  (\code{TokenMemmapDataset} in \code{llmsys.data}); use pinned memory and
  \code{non\_blocking=True} copies; keep \code{num\_workers}$>0$ if any CPU work remains.
\item \textbf{Small, unfused kernels.} LayerNorm, GELU, residual adds, dropout and the optimizer
  are memory-bound. \code{torch.compile} fuses chains of them; \code{AdamW(fused=True)} updates
  all parameters in a handful of kernels instead of several per tensor.
\item \textbf{Shapes that miss the tensor cores.} Keep $\dmodel$, $\dff$, head dimension and
  vocabulary size multiples of 64 or 128. Padding GPT-2's vocabulary from 50{,}257 to 50{,}304
  is a well-known free speed-up.
\item \textbf{Host-device synchronization.} \code{loss.item()} every step, Python-side
  \code{if} on tensor values, or printing forces the CPU to wait for the GPU and drains its
  queue of launched kernels. Log asynchronously or every $k$ steps.
\item \textbf{Attention.} Use \code{F.scaled\_dot\_product\_attention} (FlashAttention backend);
  the core models accept \code{use\_sdpa=True}.
\end{enumerate}
\end{systemslens}

% ======================================================================
\section{Scaling laws and compute budgeting}\label{sec:pretrain:scaling}

How should a fixed compute budget $C$ be split between model size $N$ and data $D$? Kaplan et
al.\ (2020) fit the loss as power laws in $N$, $D$ and $C$ and concluded that size should grow
faster than data. Hoffmann et al.\ (2022, ``Chinchilla'') revisited this with learning-rate
schedules matched to each run's length and found that $N$ and $D$ should grow \emph{equally}:
$N_{\text{opt}}\propto C^{0.5}$, $D_{\text{opt}}\propto C^{0.5}$, with roughly 20 training
tokens per parameter. They fit
\begin{equation}
L(N,D) = E + \frac{A}{N^{\alpha}} + \frac{B}{D^{\beta}},
\end{equation}
an irreducible entropy $E$ plus terms for finite model capacity and finite data. Combining
$D=20N$ with \cref{eq:pretrain:6nd} gives the planning rule
\begin{equation}
N_{\text{opt}} \approx \sqrt{\frac{C}{120}}, \qquad D_{\text{opt}}\approx 20\,N_{\text{opt}}.
\end{equation}
For $C=5.9\times10^{23}$ FLOPs this gives $N\approx70$B and $D\approx1.4$T tokens, which is
Chinchilla itself.

\begin{figure}[t]
\centering
\begin{tikzpicture}
\begin{loglogaxis}[width=0.8\linewidth, height=5.4cm, xmin=1e18, xmax=1e26,
  xlabel={training compute $C$ (FLOPs)}, ylabel={optimal size}, label style={font=\small},
  tick label style={font=\scriptsize},
  legend style={font=\scriptsize, draw=none, at={(0.03,0.97)}, anchor=north west}]
\addplot[very thick, sysblue, domain=1e18:1e26, samples=20] {sqrt(x/120)};
\addplot[very thick, napkin, domain=1e18:1e26, samples=20] {20*sqrt(x/120)};
\addplot[only marks, mark=*, mark size=1.8pt, keygreen] coordinates {(1.8e18,1.24e8) (5.88e23,7e10)};
\node[font=\scriptsize, anchor=west] at (axis cs:2.5e18,8e7) {GPT-2 small, compute-optimal};
\node[font=\scriptsize, anchor=west] at (axis cs:8e23,5e10) {Chinchilla 70B};
\legend{$N_{\text{opt}}$ (parameters), $D_{\text{opt}}$ (tokens)}
\end{loglogaxis}
\end{tikzpicture}
\caption{Chinchilla-style planning curves from $C=6ND$ and $D=20N$. Both grow as $\sqrt C$.}
\label{fig:pretrain:chinchilla}
\end{figure}

\begin{keyidea}[Compute-optimal is not deployment-optimal]
Chinchilla minimizes \emph{training} compute for a target loss. A model that will serve
billions of tokens should be smaller than compute-optimal and trained on far more tokens,
because inference cost scales with $N$. This is why recent open models are ``over-trained'':
Llama~3 8B saw about 15T tokens, nearly 2{,}000 tokens per parameter. Systems teams care about
this trade-off directly, since the serving fleet (\cref{ch:inference,ch:serving}) usually costs
more over a model's lifetime than its training run.
\end{keyidea}

\begin{napkinmath}[Planning a 7B run]
How long does a 7B model take on 2T tokens with 512 H100s at 40\% MFU? Per GPU,
$0.4\cdot989\times10^{12}/(6\cdot7\times10^9)\approx9{,}400$ tokens/s. Cluster: $4.8$M tokens/s.
$2\times10^{12}/4.8\times10^6\approx4.2\times10^5$\,s $\approx$ 4.8 days, if nothing fails. With a
cluster-wide MTBF of a few hours, checkpointing policy (\cref{sec:pretrain:save}) and
fast restarts are what decide whether it takes 5 days or 8.
\end{napkinmath}

% ======================================================================
\section{Training stability at scale}\label{sec:pretrain:stability}

Large runs fail in ways small runs do not. The most common is the \emph{loss spike}: a sudden
jump of the loss, sometimes recovering, sometimes diverging. Practical defenses, roughly in
order of how often they are used:
\begin{itemize}
\item warmup, clipping and a not-too-aggressive peak learning rate (\cref{sec:pretrain:bells});
\item lower Adam $\beta_2$ (0.95 instead of 0.999) so the second-moment estimate adapts quickly
  when gradient statistics shift;
\item pre-LayerNorm blocks (\cref{ch:gpt}), QK-normalization of queries and keys before the dot
  product (used by Qwen3, \cref{ch:modern}), and a small \emph{z-loss} $10^{-4}\cdot(\log Z)^2$ on
  the softmax normalizer to keep output logits from drifting (PaLM);
\item restarting from the last good checkpoint and skipping the data batches around the spike;
\item removing biases and using careful initialization (scaled residual initialization).
\end{itemize}

\begin{systemslens}[Determinism and debuggability]
Reproducing a spike requires bitwise-identical replay: the same data order (save the sampler
state and the data-loader position, not just the epoch), the same RNG states (Lab~7 stores
them), and deterministic kernels (\code{torch.use\_deterministic\_algorithms(True)}, which may
cost speed). At cluster scale, add per-step logging of the loss, the gradient norm per layer
group, the update-to-weight ratio $\|\Delta\theta\|/\|\theta\|$ and the largest attention logit.
These cheap scalars are the ``flight recorder'' that turns a mysterious divergence into a
diagnosable one.
\end{systemslens}

% ======================================================================
\section{Summary}\label{sec:pretrain:summary}
\begin{itemize}
\item Pretraining minimizes the average negative log-likelihood of the next token
  (cross-entropy); its exponential, perplexity, starts near the vocabulary size and is only
  comparable between models with the same tokenizer.
\item The training loop is the standard PyTorch loop; on a tiny corpus a GPT-2-sized model
  quickly memorizes the training text while the validation loss stalls.
\item Temperature and top-$k$ shape the sampling distribution; \code{temperature=0} is greedy.
\item Persist model \emph{and} optimizer \emph{and} RNG state; use safetensors for released
  weights, sharded and asynchronous checkpoints at scale, and Young's rule for the interval.
\item Loading GPT-2 requires transposing Conv1D weights, splitting the fused QKV projection,
  enabling QKV biases, and tying the output head to the token embedding.
\item Warmup, cosine (or WSD) decay and global-norm clipping make optimization robust.
\item Mixed-precision AdamW costs 16 bytes per parameter of model state; activations scale with
  $s\,b\,h\,L$ and, without FlashAttention, with $s^2$. Accumulation and checkpointing trade
  time for memory.
\item Training costs $\approx6ND$ FLOPs; MFU measures how well the hardware is used; Chinchilla
  says $D\approx20N$ for compute-optimal training, while deployment favors smaller, longer-trained
  models.
\end{itemize}

% ======================================================================
\section{Exercises}\label{sec:pretrain:exercises}

\begin{exercise}[Initial loss]
For a vocabulary of $V$ tokens and logits drawn i.i.d.\ from $\mathcal N(0,\sigma^2)$, show
that the expected initial loss is approximately $\ln V + \sigma^2/2$ for small $\sigma$. What
does this predict for GPT-2 small if the output head is initialized with standard deviation
0.02 and the final hidden states have unit variance?
\end{exercise}

\begin{exercise}[Perplexity across tokenizers]
Model A (GPT-2 tokenizer) reaches loss 3.10 on a corpus that it tokenizes into 1.00M tokens;
model B (a byte-level tokenizer) reaches loss 0.95 on the same 4.3MB of text. Convert both to
bits per byte and decide which model is better.
\end{exercise}

\begin{exercise}[Training vs.\ validation]
Run Lab~7's \code{run.py train} on a short text for 10 epochs with dropout 0.0 and with 0.1.
Plot both loss curves. At which epoch does the validation loss reach its minimum in each case,
and how do generated samples differ between those two checkpoints?
\end{exercise}

\begin{exercise}[Temperature limits]
Prove that as $\tau\to0^+$ the temperature-scaled softmax converges to a one-hot vector on the
argmax (assume a unique maximum), and that as $\tau\to\infty$ it converges to uniform. What
happens with ties at $\tau\to0$?
\end{exercise}

\begin{exercise}[Top-$k$ and temperature commute?]
Does applying temperature before or after the top-$k$ mask change the sampled distribution?
Does the same hold for top-$p$ (\cref{ch:scaling})? Justify.
\end{exercise}

\begin{exercise}[Resume correctness]
List every piece of state required for a bit-exact resume of a single-GPU run that uses
dropout, a shuffling \code{DataLoader} with two worker processes, bf16 autocast and a
\code{GradScaler}. Which of them does \cref{lst:pretrain:ckpt} miss?
\end{exercise}

\begin{exercise}[Checkpoint interval]
A job on 2{,}048 GPUs has a per-GPU MTBF of 50{,}000 hours (assume independent failures). A
synchronous checkpoint takes 90\,s; an asynchronous one stalls training for 6\,s. Compute
Young's optimal interval for both and the expected fraction of time lost to checkpointing plus
rework in each case.
\end{exercise}

\begin{exercise}[Gradient accumulation is not always exact]
Give two concrete situations in which $k$ micro-batches with accumulation do \emph{not}
reproduce the gradient of one $k$-times larger batch. (Hint: think about how the loss is
normalized when sequences contain padding, and about layers with batch statistics.)
\end{exercise}

\begin{exercise}[Activation memory]
Using \cref{eq:pretrain:act}, find the largest micro-batch size that fits GPT-2 xl
($L=48$, $h=1600$, $a=25$, $s=1024$) on an 80\,GB GPU in mixed precision, (a) without
FlashAttention, (b) with FlashAttention, (c) with full activation checkpointing. Leave 10\%
headroom.
\end{exercise}

\begin{exercise}[Measure MFU]
With \code{run.py bench} on any GPU, measure tokens/s for GPT-2 small at sequence length 1024
under: fp32; bf16 autocast; bf16 + SDPA; bf16 + SDPA + \code{torch.compile}; and the last one
with activation checkpointing. Report MFU for each and explain the ranking.
\end{exercise}

\begin{exercise}[Chinchilla planning]
You have 1{,}000 H100-days at 40\% MFU. What model size and token count are compute-optimal?
If the model must serve $10^{12}$ tokens per month for a year at $2N$ FLOPs per token, would you
rather train a model half that size on four times the data? Estimate total (training +
inference) FLOPs for both options.
\end{exercise}

\begin{exercise}[Vocabulary-chunked loss]\label{ex:pretrain:chunked}
Implement a cross-entropy that processes the sequence dimension in chunks of 128 positions,
computing logits, loss and the gradient of the hidden states for each chunk inside a custom
\code{autograd.Function}, so that the full $(B,T,V)$ logits never exist. Verify gradients
against \code{F.cross\_entropy} and measure peak memory on a GPU.
\end{exercise}

% ======================================================================
\section{Lab 7: An instrumented pretraining loop}\label{sec:pretrain:lab}

\begin{labbox}{7}{An instrumented pretraining loop}
\labdir{lab07_pretrain}

\labsection{Goals}
Build, from the loss up, a pretraining loop you would trust with real compute: correct
objective, robust schedule, exact resume, memory-saving switches, and the napkin-math
instruments to judge its efficiency.

\labsection{Setup}
\code{cd code \&\& pip install -r requirements.txt}. All tests run on CPU in under a minute
with a tiny GPT and a synthetic, noisy periodic token stream:
\code{LAB\_IMPL=starter python -m pytest -q labs/lab07\_pretrain}.

\labsection{Tasks}
\begin{enumerate}
\item \ident{cross_entropy_from_scratch} and \ident{perplexity}: log-sum-exp formulation; must
  match \code{F.cross\_entropy} and stay finite for logits of $\pm1000$.
\item \ident{warmup_cosine_lr}: linear warmup from \code{initial\_lr} to \code{peak\_lr}, then
  cosine decay to \code{min\_lr} (\cref{eq:pretrain:cosine}).
\item \ident{global_grad_norm} and \ident{clip_grad_norm_}: must agree with
  \code{torch.nn.utils.clip\_grad\_norm\_}.
\item \ident{generate_text}: greedy, temperature, top-$k$, EOS, seeded generator.
\item \ident{save_checkpoint} / \ident{load_checkpoint}: 4+4 steps with a restore in between
  must equal 8 straight steps bit-for-bit.
\item \ident{enable_activation_checkpointing}: identical gradients to the plain model.
\item \ident{train_model}: warmup+cosine, clipping, accumulation, bf16 autocast, periodic
  evaluation, logging of LR and pre-clip gradient norm, checkpoints.
\item \ident{training_memory_estimate}, \ident{chinchilla_optimal}, \ident{mfu}.
\item \ident{load_openai_gpt2_params}: OpenAI's original nested layout (transposes, QKV split,
  tying).
\end{enumerate}

\labsection{Measurements}
\begin{itemize}
\item \code{python run.py train --text the-verdict.txt --epochs 10}: plot train/val loss;
  you should reproduce the overfitting picture of \cref{sec:pretrain:train} (val loss bottoms
  out around epoch 2--4 while train loss keeps falling).
\item Repeat with a 10$\times$ larger text (any public-domain novel) and compare the gap.
\item On a GPU: \code{python run.py bench --preset 124M --seq 1024} with and without
  \code{--amp bf16}, \code{--sdpa}, \code{--compile}, \code{--ckpt-act}. Record ms/step, peak
  memory and MFU; compare peak memory with \ident{training_memory_estimate}. Expect bf16+SDPA
  to give the largest jump, compile a further gain, and checkpointing to cut activation memory
  by an order of magnitude at roughly 25--35\% lower throughput.
\item \code{python run.py sample --hf-dir ./gpt2}: sample from the real GPT-2 at temperatures
  0, 0.7 and 1.5 with top-$k$ 25.
\end{itemize}

\labsection{Deliverables}
Passing tests; a one-page report with the loss curves, the benchmark table (with MFU) and a
comparison of measured vs.\ estimated memory.

\labsection{Stretch goals}
Add a \code{GradScaler} path for fp16 and show a run that diverges without it; implement the
chunked cross-entropy of \cref{ex:pretrain:chunked}; add a WSD schedule and
branch two annealing runs from the same stable-phase checkpoint.
\end{labbox}
